% Zdroj: PDF priklady-jaro-2026.pdf, fyzická str. 21. Značka: specializace UI-SU, UI-ZPJ. Zařazeno: S1 (CSP / hranová konzistence).
\section{CSP}
\szzotazka{jaro 2026}{27}{CSP (binární CSP, hranová konzistence, AC-3)}

\noindent\textbf{Zadání.}\par
\begin{enumerate}[nosep]
  \item Definujte binární Constraint Satisfaction Problem (CSP), konkrétně vstup a výstup tohoto problému. Co znamená ,,hranová konzistence`` v kontextu CSP?
  \item Pokud je problém hranově konzistentní, je zaručeno, že je splnitelný? Pokud problém není hranově konzistentní, může být stále splnitelný? Svou odpověď zdůvodněte nebo poskytněte protipříklady.
  \item Uvažte následující binární CSP, aplikujte algoritmus AC-3 pro dosažení hranové konzistence. Součástí řešení musí být postup.

  Mějme celočíselné proměnné $A, B, C$ a $D$ s doménami:
  \begin{align*}
    D_A &= \{1\} \\
    D_B &= \{1, 2, 3\} \\
    D_C &= \{1, 2, 3\} \\
    D_D &= \{1, 2, 3\}
  \end{align*}
  Podmínky:
  \begin{align*}
    C_1 &: C + D = 4 \\
    C_2 &: A + B \leq 3 \\
    C_3 &: B + C \geq 4
  \end{align*}
\end{enumerate}

\medskip
\noindent\textbf{Náčrt řešení.}\par
\begin{enumerate}[nosep]
  \item CSP je definován množinou proměnných, doménou (tj. množinou hodnot) pro každou proměnnou a množinou podmínek nad proměnnými. V binárním CSP může každá podmínka obsahovat maximálně dvě proměnné. Podmínka je hranově konzistentní, pokud pro každou hodnotu z domény každé proměnné v dané podmínce existuje přípustná hodnota z domény ostatních proměnných z dané podmínky. Problém je hranově konzistentní, pokud je každá podmínka hranově konzistentní.
  \item Hranově konzistentní problém nemusí být splnitelný. Například tři proměnné $X_1, X_2, X_3$ s binárními doménami $\{0, 1\}$ a podmínkami pro každou dvojici $X_i \neq X_j$. Hranově nekonzistentní problém může být splnitelný. Například příklad ze zadání.
  \item Algoritmus AC-3 mění domény proměnných tak, aby byly všechny podmínky hranově konzistentní. Podmínky, které obsahují proměnné, kterým se změnily domény, se musí zkontrolovat znovu. Hranově konzistentní problém má domény $D_A = \{1\}$, $D_B = \{1, 2\}$, $D_C = \{2, 3\}$, $D_D = \{1, 2\}$.
\end{enumerate}

\medskip
\noindent\textit{Doplňující vysvětlení (nad rámec oficiálního náčrtu):} Postup AC-3 (revize hrany $(X,Y)$ vyřadí z~$D_X$ hodnoty bez opory v~$D_Y$; po zúžení $D_X$ se do fronty vrátí hrany $(Z,X)$ od ostatních sousedů $Z$):
\begin{enumerate}[nosep]
  \item $C_2$, revize $(B,A)$: $A=1$ a~$A+B\le3$ dává $B\le2$, vyřadí se $B=3$, $D_B=\{1,2\}$. Revize $(A,B)$ nic nemění.
  \item $C_3$, revize $(C,B)$: $B\le2$ a~$B+C\ge4$ dává $C\ge2$, vyřadí se $C=1$, $D_C=\{2,3\}$. Revize $(B,C)$ nic nemění ($b=1$ i~$b=2$ má oporu $c=3$).
  \item Zúžila se $D_C$, do fronty se vrací $(D,C)$: $C+D=4$ a~$C\in\{2,3\}$ dává $D\in\{1,2\}$, vyřadí se $D=3$.
  \item Zúžila se $D_D$, do fronty se vrací $(C,D)$: $c=2$ má oporu $d=2$, $c=3$ oporu $d=1$, beze změny. Fronta se vyprázdní.
\end{enumerate}
Výsledek odpovídá náčrtu. Protože původní domény obsahovaly hodnoty bez opory (např.\ $B=3$), zadaný problém hranově konzistentní nebyl, a~přesto je splnitelný (např.\ $A=1$, $B=2$, $C=2$, $D=2$); to je protipříklad, na který náčrt v~bodě~2 odkazuje.
